% Problem record op_90a05e57e44c086e. This TeX file is derived from
% database/problems_json/op_90a05e57e44c086e.json by scripts/sync-tex.mjs; edit the JSON record
% and rerun the script rather than editing this file.
\section{Polynomial-time quantum algorithm for the Dihedral Hidden Subgroup Problem}
\label{sec:op_90a05e57e44c086e}

\paragraph{Problem.}
Does the Dihedral Hidden Subgroup Problem admit a quantum algorithm whose
running time is polynomial in the input length?

For a positive integer $N$, let the dihedral group be
\begin{equation}
  D_N=\langle r,s\mid r^N=1,\ s^2=1,\ srs=r^{-1}\rangle ,
  \label{eq:90a0-dihedral-presentation}
\end{equation}
so that $D_N$ has order $2N$.  Consider an oracle $f:D_N\to X$ promised to
hide a subgroup generated by an unknown reflection.  Equivalently, for an
unknown $a\in\mathbb Z_N$, let
\begin{equation}
  H_a=\{1,sr^a\},
  \qquad
  f(x)=f(y)\iff xH_a=yH_a .
  \label{eq:90a0-reflection-subgroup}
\end{equation}
Given quantum oracle access to $f$, the task is to recover $a$, and hence
$H_a$ in \eqref{eq:90a0-reflection-subgroup}.  The open question is whether
this can be done with bounded error using a number of elementary quantum
operations polynomial in $\log N$, where $D_N$ is defined by
\eqref{eq:90a0-dihedral-presentation}.

\subsection*{Status}

Unsolved

\subsection*{Source}

Implicit in the Dihedral Hidden Subgroup Problem literature, with
the current-status formulation supported by Chen and Sun
\sourcecite{ref:90a0-chen-sun}{CS24}.  The wording above is a contributor
formulation rather than a verbatim open question attributed to Kuperberg,
Regev, or Chen and Sun.

\subsection*{Progress}

\begin{itemize}
  \item Kuperberg developed subexponential-time quantum algorithms for the
Dihedral Hidden Subgroup Problem; his collimation-sieve formulation runs in
$\exp(O(\sqrt{\log N}))$ quantum time, uses $O(\log N)$ quantum space and
subexponential classical space, and no polynomial-time algorithm follows
from this sieve framework
\sourcecite{ref:90a0-kuperberg}{Kup13}.

  \item Regev showed that an efficient solution to the relevant dihedral coset
problem would yield a quantum algorithm for certain instances of unique
shortest vector problems, establishing a direct link between the dihedral
problem and quantum algorithms for lattice problems
\sourcecite{ref:90a0-regev}{Reg04}.

  \item Chen and Sun give a modern account of the Dihedral Hidden Subgroup
Problem, including obstructions to standard Fourier sampling and bounds for
algorithms acting on dihedral coset states; their survey treats the
quantum complexity of the problem as unresolved and records subexponential
algorithms as the established general algorithmic regime
\sourcecite{ref:90a0-chen-sun}{CS24}.

  \item A polynomial-time algorithm for the closely related Dihedral Coset
Problem was claimed in a 2026 preprint, but Gupte, Ragavan, and Zhandry
proved that the proposed algorithm cannot recover even the least-significant
bit of the secret with non-negligible advantage and therefore does not solve
the Dihedral Coset Problem; their no-go result also indicates that
successful Regev-type constructions may need to retain and exploit
substantially more classical Fourier-label information during uncomputation
\sourcecite{ref:90a0-grz}{GRZ26}.
\end{itemize}

\subsection*{References}

\begin{enumerate}[label={},leftmargin=4.8em,labelsep=0.5em]
  \footnotesize
  \item[\textup{[Kup13]}]\label{ref:90a0-kuperberg}
  Greg Kuperberg, ``Another Subexponential-time Quantum Algorithm for the
Dihedral Hidden Subgroup Problem,'' in \emph{8th Conference on the Theory of
Quantum Computation, Communication and Cryptography (TQC 2013)}, LIPIcs
\textbf{22}, 20--34 (2013).
\href{https://doi.org/10.4230/LIPIcs.TQC.2013.20}{doi:10.4230/LIPIcs.TQC.2013.20};
\href{https://arxiv.org/abs/1112.3333}{arXiv:1112.3333}.

  \item[\textup{[Reg04]}]\label{ref:90a0-regev}
  Oded Regev, ``Quantum Computation and Lattice Problems,'' \emph{SIAM
Journal on Computing} \textbf{33}, 738--760 (2004).
\href{https://doi.org/10.1137/S0097539703440678}{doi:10.1137/S0097539703440678};
\href{https://arxiv.org/abs/cs/0304005}{arXiv:cs/0304005}.

  \item[\textup{[CS24]}]\label{ref:90a0-chen-sun}
  Imin Chen and David Sun, ``The Dihedral Hidden Subgroup Problem,'' \emph{Journal
of Mathematical Cryptology} \textbf{18}, Article 20220029 (2024).
\href{https://doi.org/10.1515/jmc-2022-0029}{doi:10.1515/jmc-2022-0029};
\href{https://arxiv.org/abs/2106.09907}{arXiv:2106.09907}.

  \item[\textup{[GRZ26]}]\label{ref:90a0-grz}
  Aparna Gupte, Seyoon Ragavan, and Mark Zhandry, ``The ePrint:
2026/1591 Quantum Algorithm Does Not Solve DCP,'' Cryptology ePrint
Archive, Paper 2026/1693 (2026).
\href{https://eprint.iacr.org/2026/1693}{ePrint 2026/1693}.
\end{enumerate}

\subsection*{Comment}

The Dihedral Hidden Subgroup Problem is one of the central instances of
non-Abelian hidden subgroup problems and, through Regev's reduction (Progress
above), is directly tied to quantum algorithms for lattice problems; see the related problem on \href{https://qiqc-op.com/problem/op_3596f8c955c66d94/}{a polynomial-time quantum algorithm for Learning With Errors}.  No polynomial-time quantum algorithm is known
in the standard oracle model.

\subsection*{Field}

Quantum algorithm

\subsection*{Topic}

Computational complexity and computability

\subsection*{ID}

\texttt{op\_90a05e57e44c086e}
